% Junior Level 2, Week 10 -- answer key
\answersweek{10}{Arcade Showcase}{Answers for workbook chapter 10 and play-at-home card 10.}

\section{Lesson questions}
\begin{answer}[Exit question]
The game would end after 5 misses instead of 3, so it would be easier: the player can miss more apples (or bananas) before it is game over. Any answer that links the number to how easy or hard the game is shows the child understands the design choice.
\end{answer}

\begin{answer}[Reflection]
Any drawing of a moment from the arcade and any honest smiley is right. Write down the answers to \enquote{What would you like to make next?}: they help you plan Junior Level 3 (Bridge to Python), where children plan bigger programs and meet Python.
\end{answer}

\begin{answer}[Show-me checklist]
The Week 10 show-me checklist is in the Assessments booklet. It is observed during the showcase and the show-me corner and is never graded. Compare it with the week 5 checklist: most children move at least one row up. \enquote{Getting there} is a normal result at this age.
\end{answer}

\section{Workbook exercises}
\begin{answer}[Exercise 10.1: Last-Minute Bug Check]
A finished game gets \enquote{yes} in all four rows. Common \enquote{not yet} answers: row 1 or 3, the second game starts with the old score or says Game over at once (fix: \enquote{set score to 0} and \enquote{set misses to 0} at the top of the referee script); row 2, the game never stops (fix: \enquote{stop all} inside the if, and the if inside the forever); row 4, no sound (check the computer volume first, then the sound block). A sprite that stays hidden on the second game needs \enquote{show} under its green flag. Fixing one thing is enough today.
\end{answer}

\begin{answer}[Exercise 10.2: My One-Minute Demo]
Any true sentences about the child's own game are right, for example \enquote{My game is called Ghost Maze. You are the Cat. You win when you reach the apple. I chose a small cat, because the big one got stuck in the walls.} Look for: the name, the player, how to win, and one design choice with a reason. A reason that came from a Player's wish is the best kind. Pointing at the script that makes the choice happen counts as \enquote{Got it} for explaining.
\end{answer}

\begin{answer}[Exercise 10.3: Arcade Stickers]
Any sticker placed with a kind spoken reason is right, plus two stars and a wish (\enquote{I liked the pop and the ghost. I wish it was a bit slower.}). Look for: the visitor asks before playing, uses kind words only, plays to the end and does not change the game.
\end{answer}

\section{Play-at-home card}
\begin{answer}[Family arcade]
Look for the same demo as in class: how to win and lose, the family member plays without help, and the child listens to the stars and the wish. The child does not need to change the game afterwards.
\end{answer}

\begin{answer}[What next?]
Any idea is right. Second levels, mazes, two-player games and drawing games all lead into Junior Level 3.
\end{answer}
